\documentclass{article}

% Language setting
\usepackage[english]{babel}

% Set page size and margins
\usepackage[letterpaper,top=2cm,bottom=2cm,left=3cm,right=3cm,marginparwidth=1.75cm]{geometry}

% Useful packages
\usepackage[utf8]{inputenc}
\usepackage{latexsym}
\usepackage{amsfonts}
\usepackage{amssymb}
\usepackage{amsthm}
\usepackage{amsmath}
\usepackage{tikz}
\usepackage{circuitikz} 
\usepackage{listings}
\usepackage{mathtools}
\usepackage{float}
\usepackage{natbib}
\usepackage{fourier}
\bibpunct{[}{]}{,}{n}{}{,~}
\usepackage{minted}

\title{Software for Data Science: Julia}
\date{24/11/2025}

\begin{document}
\maketitle
\section*{Exercise 1}

\medskip
\noindent
Find bellow a list of topics covered in this Julia lecture. Choose 4 topics. Take some time to go over the lecture again. Prepare a presentation of these topics for the group.

\begin{itemize}
    \item Packages and working environment;
    \item Comments, printing;
    \item Comparison and arithmetic operations;
    \item Control flows (if, while, for, break, continue);
    \item Data (dictionary, dataframes);
    \item Functions (keyword arguments, optional arguments, bang convention, broadcasting);
    \item Error handling;
    \item Scoping and closures;
    \item Modules;
    \item Typing (abstract, concrete, conversion, promotion);
    \item User defined types (composite types, mutability, constructors);
    \item Basic hardware;
    \item Type stability.
\end{itemize}

\section*{Exercise 2}

\medskip
\noindent
\texttt{DataFrames.jl} provides in‑memory tabular data structures similar to R \texttt{data.frames} or Python’s pandas \texttt{DataFrame}.

\begin{minted}{julia}
using DataFrames
\end{minted}

\medskip
\noindent
A \texttt{DataFrame} is a table: columns have names and types, rows are observations.
\begin{itemize}
    \item Columns are vectors of equal length; different columns can have different types (e.g., Int64, Float64, String, Union{Missing, T}).
    \item Missing values are represented with missing and handled via Union{Missing, T} column types.
\end{itemize}

\medskip
\noindent
Even though in pratice data will usually not be added by hand in the \texttt{DataFrame}, it can be done the following way.

\begin{minted}{julia}
df = DataFrame(
    id      = 1:2,
    name    = ["Alice", "Thomas"],
    score   = [8.5, 7.0],
    passed  = [true, true]
)
\end{minted}

\paragraph{Import and clean}

\medskip
\noindent
In this exercise, we will be working with the following data.

\begin{minted}{julia}
df = DataFrame(
    id      = 1:10,
    name    = ["Alice", "Thomas", "Xiaobo", "Imane", "Lucia", "Rohan", "Camille", 
    "Sophie", "Fatoumata", "Noé"],
    score   = [-1, 13.5, 14, 12, 16, missing, 10, 15, 17, 7],
    passed  = [false, true, true, true, true, false, true, true, true, false]
)
\end{minted}

\medskip
\noindent
Write this table to a CSV file using \texttt{write} from \texttt{CSV.jl}. Then load this file using the \texttt{read} method.

\medskip
\noindent
Have a close look at the data in the table. Use \texttt{filter} in order to sort out missing students and invalid grades, i.e. each grade $x$ should $0 \ge x \ge 20$. To check for the value \texttt{missing} use \texttt{ismissing}. 

\paragraph{Exploratory analysis}

\medskip
\noindent
Use the \texttt{mean}, \texttt{minimum} and \texttt{minimum} to compute statistics on the students grades. Look into the documentation of Julia to find the function to determine the minimum/maximum of an array and the associated index. Which student go the best/worse grade ?

\paragraph{Relational analysis}

\medskip
\noindent
Let the bellow table be the gender associated to the students. Operate a \texttt{join} between \texttt{df} and \texttt{df\_gender} to join gender information with exam records. What is the difference between the \texttt{innerjoin}, \texttt{outerjoin} and \texttt{leftjoin} functions ?

\begin{minted}{julia}
df_gender = DataFrame(
    id      = 1:10,
    gender  = ["F", "M", "M", "F", "F", "M", "M", "F", "F", "M"],
)
\end{minted}

\paragraph{Performance}

\medskip
\noindent
Compare the performance between row-wise looping and column-wise transformations using \texttt{BenchmarkTools.jl}. A standard \texttt{for} loop will be used to go over the rows and to act on the column broadcasting will be used.

\begin{minted}{julia}
N = 1000000
df = DataFrame(value = rand(N))
\end{minted}

\medskip
\noindent
To do this comparison, multiply each value in the data set by 2. Evaluate the time required for this operation on the whole table with \texttt{@btime}. 

\end{document}